%%%%%%%%%%%%%%%%%
% CV tailored for engIT - Ingénieur Machine Learning embarqué
% Positioning: Ingénieur Machine Learning / Performance | Python, C/C++, Linux, validation, TensorFlow/ONNX ramp-up & systèmes contraints
%%%%%%%%%%%%%%%%%

\documentclass[8pt,a4paper,ragged2e,withhyper]{altacv}

\geometry{left=1.25cm,right=1.25cm,top=.5cm,bottom=1.cm,columnsep=1.2cm}

\usepackage{paracol}
\usepackage{fontawesome5}
\usepackage{hyperref}

\iftutex
  \setmainfont{Roboto Slab}
  \setsansfont{Lato}
  \renewcommand{\familydefault}{\sfdefault}
\else
  \usepackage[rm]{roboto}
  \usepackage[defaultsans]{lato}
  \renewcommand{\familydefault}{\sfdefault}
\fi

\definecolor{SlateGrey}{HTML}{2E2E2E}
\definecolor{LightGrey}{HTML}{666666}
\definecolor{DarkPastelRed}{HTML}{8AC0D9}
\definecolor{PastelRed}{HTML}{00B0DA}
\definecolor{GoldenEarth}{HTML}{B4AD0C}

\colorlet{name}{black}
\colorlet{tagline}{PastelRed}
\colorlet{heading}{DarkPastelRed}
\colorlet{headingrule}{GoldenEarth}
\colorlet{subheading}{PastelRed}
\colorlet{accent}{PastelRed}
\colorlet{emphasis}{SlateGrey}
\colorlet{body}{LightGrey}

\definecolor{violet}{HTML}{5A2582}
\definecolor{rouge}{HTML}{F53B3B}
\definecolor{noir}{HTML}{00232C}
\definecolor{bleu}{HTML}{00B0DA}
\definecolor{rose}{HTML}{F831A0}
\definecolor{jaune}{HTML}{FFEC00}
\definecolor{vert}{HTML}{34AD6C}

\renewcommand{\namefont}{\Huge\rmfamily\bfseries}
\renewcommand{\personalinfofont}{\footnotesize}
\renewcommand{\cvsectionfont}{\LARGE\rmfamily\bfseries}
\renewcommand{\cvsubsectionfont}{\large\bfseries}
\renewcommand{\cvItemMarker}{{\small\textbullet}}
\renewcommand{\cvRatingMarker}{\faCircle}

\input{../../_shared/pubs-num.tex}

% auto_tex_manager layout tuning
\emergencystretch=3em
\hfuzz=60pt
\hbadness=9999
\sloppy

\begin{document}

\name{BOILEAU Guillaume}
\tagline{Ingénieur Machine Learning / Performance | Python, C/C++, Linux, validation \& systèmes contraints}

\photoL{2.8cm}{../../_shared/moi}

\personalinfo{%
  \email{guillaume.boileaupro@gmail.com}
  \phone{+33 6 59 94 85 84}
  \location{Alpes-Maritimes, FRANCE}
  \linkedin{guillaume-boileau-891b81265}
  \researchgate{Guillaume-Boileau-2}
  \orcid{0000-0002-3576-6968}
}

\makecvheader
\vspace{-0.3cm}

\AtBeginEnvironment{itemize}{\small}
\columnratio{0.64}

\begin{paracol}{2}

\cvsection{Experience}

\cvevent{\textbf{Software Engineer -- Développement logiciel, intégration \& support opérationnel}}{VEV Platform Services France}{2025 -- 2026}{Valbonne, France}

\textbf{Contexte :} Plateforme logicielle industrielle CPMS connectée à des infrastructures de recharge, avec services applicatifs, données opérationnelles, flux d’échanges, contraintes terrain et besoins d’exploitation.

\textbf{Objectif :} Contribuer au développement, à l’intégration et à la fiabilisation de services logiciels opérationnels, en combinant analyse de flux, investigation d’incidents, validation de comportements applicatifs et documentation technique.

\begin{itemize}
  \item \textbf{Développement logiciel :} Contribution à des composants TypeScript, Angular et Node.js intégrés dans une plateforme opérationnelle.
  \item \textbf{Intégration logiciel / terrain :} Travail sur les interfaces entre services applicatifs, données opérationnelles, équipements connectés et contraintes d’exploitation.
  \item \textbf{Analyse d’incidents :} Investigation d’échanges FTP, interruptions de flux, incohérences de données et comportements inattendus.
  \item \textbf{Validation opérationnelle :} Vérification de la cohérence entre données applicatives, comportement des équipements et attendus fonctionnels.
  \item \textbf{Support technique :} Diagnostic, analyse de causes, formalisation de constats et contribution aux actions de correction.
  \item \textbf{Pratiques projet :} Utilisation de Git, Linux, Zenhub, documentation technique et suivi collaboratif des sujets.
\end{itemize}

\divider

\cvevent{\textbf{Ingénieur R\&D senior -- Machine learning, simulation \& validation de performance}}{CNES}{2023 -- 2025}{Nice, France}

\textbf{Contexte :} Développement logiciel et R\&D pour le segment sol de la mission spatiale LISA ESA/NASA, avec chaînes de données mission L0--L3, simulation, intégration de modèles, analyse de performance et validation technique en environnement CNES/ESA.

\textbf{Objectif :} Développer des workflows Python robustes pour intégrer, comparer, valider et exploiter des modèles et données complexes, avec contraintes de performance, reproductibilité et fiabilité des résultats.

\begin{itemize}
  \item \textbf{Python scientifique \& logiciel :} Développement de CATpyLISA, plateforme Python pour intégration, comparaison, validation et revue de modèles.
  \textcolor{DarkPastelRed}{\href{https://gitlab.oca.eu/gboileau/lisa_l3_tools}{\bf GitLab: CATpyLISA}}
  \item \textbf{Machine learning / data-driven workflows :} Structuration de pipelines d’analyse, validation de modèles, métriques, contrôles de cohérence et comparaison de configurations.
  \item \textbf{Performance \& systèmes contraints :} Analyse de traitements numériques, volumes de données, comportements inattendus, limites de performance et effets de configuration.
  \item \textbf{Validation \& tests :} Définition de critères d’acceptation, règles de comparaison, procédures de vérification et analyse d’écarts.
  \item \textbf{Simulation \& prototypage :} Développement d’environnements numériques permettant de tester des modèles et d’évaluer leur robustesse dans différents scénarios.
  \item \textbf{Documentation technique :} Formalisation des hypothèses, méthodes, résultats, limites et éléments de revue.
  \item \textbf{Coordination technique :} Interface avec contributeurs scientifiques, logiciels et institutionnels dans un contexte international.
\end{itemize}

\divider

\cvevent{\textbf{Data Scientist -- Machine learning, capteurs \& réduction de bruit}}{Universiteit Antwerpen, Belgium}{2021 -- 2023}{Antwerp, Belgium}

\textbf{Contexte :} Travaux de data science pour instruments de précision et détecteurs de nouvelle génération, combinant données capteurs, bruit environnemental, modélisation statistique, machine learning, deep learning et validation de performance.

\textbf{Objectif :} Développer des workflows robustes pour analyser des données capteurs, réduire des bruits non stationnaires, comparer des configurations et améliorer la performance de systèmes complexes.

\begin{itemize}
  \item \textbf{Machine learning appliqué :} Structuration d’approches de réduction de bruits non stationnaires avec canaux témoins, machine learning et deep learning.
  \item \textbf{Données capteurs :} Analyse de mesures environnementales, propagation, corrélations spatiales et impacts sur la performance système.
  \item \textbf{Pipelines Python :} Développement de workflows pour analyse de données, statistiques, performance, validation et reporting.
  \item \textbf{Modélisation avancée :} Mise en place de pipelines MCMC pour différentes configurations instrumentales et différents niveaux de bruit.
  \item \textbf{Évaluation de performance :} Figures de mérite, analyses de sensibilité, contrôles de cohérence et comparaison de configurations.
  \item \textbf{Systèmes contraints :} Analyse de signaux très faibles dans des environnements bruités, avec contraintes de sensibilité, robustesse et qualité de mesure.
\end{itemize}

\divider

\cvevent{\textbf{Ingénieur R\&D junior -- Signal, simulation \& calcul scientifique}}{ARTEMIS, Observatoire de la Côte d’Azur}{2018 -- 2021}{Nice, France}

\textbf{Contexte :} Recherche appliquée liée à l’analyse de signaux faibles, à la simulation numérique, au traitement du signal et à la préparation de missions spatiales.

\textbf{Objectif :} Développer des méthodes logicielles d’analyse et de simulation afin de produire des résultats scientifiques robustes, documentés et exploitables.

\begin{itemize}
  \item \textbf{Traitement du signal :} Méthodes pour analyser et séparer des signaux faibles et superposés.
  \item \textbf{Simulation numérique :} Développement d’approches de simulation pour environnements complexes et grands jeux de données.
  \item \textbf{Calcul scientifique :} Workflows Python/C d’analyse, visualisation, comparaison et validation de résultats.
  \item \textbf{Validation :} Vérification de cohérence, comparaison de résultats, études de sensibilité et analyse de limites méthodologiques.
  \item \textbf{Documentation :} Formalisation des méthodes, hypothèses, résultats et conclusions pour revue collaborative.
\end{itemize}

\switchcolumn

\cvsection{Profile}

\textbf{Ingénieur docteur orienté Machine Learning, performance et validation de systèmes complexes, avec une expérience solide en Python, C/C++, Linux, data science, traitement du signal, simulation, tests et amélioration continue.}

Positionnement pour ce rôle : contribuer au développement et à l’optimisation de modèles de Machine Learning sur systèmes contraints, avec une montée ciblée sur TensorFlow Lite, ONNX, Linux embarqué, ARM et accélérateurs IA.

\begin{itemize}
  \item Solide socle ML/data science : apprentissage supervisé, validation de modèles, métriques, analyse d’erreurs, robustesse et performance.
  \item Expérience en données capteurs, réduction de bruit, signaux faibles, simulation, contraintes de sensibilité et systèmes complexes.
  \item Python avancé, C/C++, Linux, Git, CI/CD, workflows reproductibles et documentation technique.
  \item Capacité à collaborer avec équipes software, data, système et hardware grâce à une culture interdisciplinaire.
  \item Montée rapide visée sur ML embarqué : TensorFlow Lite, ONNX, Linux embarqué, ARM, optimisation mémoire/consommation.
\end{itemize}

\cvsection{Languages}

\cvskill{English}{4}
\divider
\cvskill{Spanish}{2}
\divider

\medskip

\cvsection{Education}

\cvevent{PhD in Physics}{Université Côte d’Azur}{2018 -- 2021}{Nice, France}

\divider

\cvevent{Selective Graduate Program in Fundamental Physics (Magistère)}{Université Paris-Saclay}{2016 -- 2018}{Orsay, France}

\divider

\cvevent{MSc in Astronomy and Astrophysics}{Observatoire de Paris / Université Paris-Saclay}{2017 -- 2018}{Paris, France}

\divider

\cvevent{BSc in Fundamental Physics}{Université Paris-Saclay}{2015 -- 2016}{Orsay, France}

\cvsection{Professional Training}

\cvevent{Machine Learning in Python with scikit-learn}{FUN MOOC -- Inria}{2026}{France Université Numérique}

\small{
Predictive modelling, supervised learning, model evaluation, cross-validation, metrics, pipelines and practical machine learning workflows with Python and scikit-learn.
}

\divider

\cvevent{Machine Learning in Python with scikit-learn}{FUN MOOC -- Inria}{2026}{Achievement Badge}

\small{
Predictive modelling with scikit-learn and methodological/software fundamentals of machine learning in Python.
}

\divider

\cvevent{Supervised Machine Learning: Regression \& Classification}{DeepLearning.AI -- Andrew Ng}{2020}{Coursera}

\divider

\cvevent{Continuous Professional Training -- Software, Data, AI and Systems}{OpenClassrooms}{2024 -- 2026}{}

\small{
Python, SQL, Machine Learning, AI, Linux, JavaScript, TypeScript, Angular, Node.js, Django, REST APIs, C/C++, algorithms and software engineering fundamentals.
}

\end{paracol}

\cvsection{Projects}

\cvevent{Noise reduction and predictive modelling workflows}{Machine learning / Sensor data / Precision instrumentation}{2021 -- 2023}{}

Développement de workflows numériques et statistiques pour la réduction de bruit non stationnaire, l’analyse de données capteurs, la comparaison de configurations et la validation de performance sur systèmes complexes.

\textbf{Rôle :} analyse de données, modélisation statistique, pipelines MCMC, approches ML/deep learning pour la réduction de bruit, validation de performance et reporting technique.

\textbf{Compétences mobilisées :} Python, machine learning, deep learning concepts, données capteurs, traitement du signal, statistiques, validation, analyse de performance.

\divider

\cvevent{CATpyLISA -- Plateforme Python de simulation, validation \& performance}{Mission spatiale LISA ESA/NASA -- CNES/ESA}{2023 -- 2025}{}

Plateforme Python dédiée à l’intégration de modèles, à la structuration de données, à la comparaison de résultats, à la validation technique et à la revue collaborative de workflows scientifiques.

\textbf{Rôle :} développement Python, structuration de workflows, simulation, validation de modèles, analyse d’écarts, documentation technique et coordination avec plusieurs contributeurs.

\textbf{Compétences mobilisées :} Python, NumPy, SciPy, pandas, Matplotlib, GitLab, simulation, validation, performance numérique, data quality, documentation.

\divider

\cvevent{VEV -- Logiciel opérationnel, intégration \& support}{Industrial software / Data flows / Connected infrastructure}{2025 -- 2026}{}

Développement et intégration de services logiciels pour une plateforme CPMS connectée à des infrastructures de recharge, avec analyse de flux FTP, investigation d’incidents, vérification de comportements applicatifs et support opérationnel.

\textbf{Rôle :} développement TypeScript, Angular et Node.js, intégration logicielle, analyse de flux, investigation d’incidents, validation opérationnelle et documentation technique.

\textbf{Compétences mobilisées :} TypeScript, Angular, Node.js, Linux, Git, analyse d’incidents, intégration logicielle, support opérationnel.

\cvsection{Compétences clés}

\vspace{-.3cm}

\cvsubsection{Machine Learning \& Deep Learning}

{\small
\cvtag{Machine Learning}
\cvtag{Deep Learning concepts}
\cvtag{scikit-learn}
\cvtag{TensorFlow -- awareness}
\cvtag{TensorFlow Lite -- ramp-up}
\cvtag{ONNX -- ramp-up}
\cvtag{Validation de modèles}
\cvtag{Métriques}
\cvtag{Analyse d’erreurs}
\cvtag{Robustesse}
\cvtag{Optimisation de modèles -- ramp-up}
\cvtag{Déploiement ML embarqué -- ramp-up}
}

\divider\smallskip
\vspace{-.3cm}

\cvsubsection{Systèmes contraints \& embarqué}

{\small
\cvtag{Systèmes contraints}
\cvtag{Linux}
\cvtag{Linux embarqué -- ramp-up}
\cvtag{ARM -- awareness}
\cvtag{Accélérateurs IA -- awareness}
\cvtag{Optimisation performance}
\cvtag{Consommation / mémoire -- awareness}
\cvtag{Temps de calcul}
\cvtag{Capteurs}
\cvtag{Instrumentation}
\cvtag{Software / hardware interface -- awareness}
}

\divider\smallskip
\vspace{-.3cm}

\cvsubsection{Python, C/C++ \& outils}

{\small
\cvtag{Python}
\cvtag{C}
\cvtag{C++}
\cvtag{Bash}
\cvtag{NumPy}
\cvtag{SciPy}
\cvtag{pandas}
\cvtag{Matplotlib}
\cvtag{Jupyter}
\cvtag{Git}
\cvtag{GitLab}
\cvtag{CI/CD}
\cvtag{Docker}
\cvtag{Linux}
}

\divider\smallskip
\vspace{-.3cm}

\cvsubsection{Validation, tests \& performance}

{\small
\cvtag{Tests}
\cvtag{Validation}
\cvtag{Amélioration continue}
\cvtag{Analyse de performance}
\cvtag{Benchmarking}
\cvtag{Contrôles de cohérence}
\cvtag{Analyse de sensibilité}
\cvtag{Diagnostic technique}
\cvtag{Root cause analysis}
\cvtag{Documentation technique}
\cvtag{Reporting}
}

\divider\smallskip
\vspace{-.3cm}

\cvsubsection{Data, signal \& simulation}

{\small
\cvtag{Data Science}
\cvtag{Données capteurs}
\cvtag{Traitement du signal}
\cvtag{Simulation numérique}
\cvtag{Analyse de bruit}
\cvtag{Signaux faibles}
\cvtag{Statistiques}
\cvtag{MCMC}
\cvtag{Modélisation}
\cvtag{Data quality}
\cvtag{Reproductibilité}
}

\divider\smallskip
\vspace{-.3cm}

\cvsubsection{Collaboration \& environnement international}

{\small
\cvtag{Équipe multidisciplinaire}
\cvtag{Software / data / hardware}
\cvtag{Coordination technique}
\cvtag{Documentation}
\cvtag{Confluence}
\cvtag{Jira}
\cvtag{Zenhub}
\cvtag{Anglais technique}
\cvtag{Environnement international}
\cvtag{Autonomie}
\cvtag{Rigueur}
}

\end{document}
